\section{Ghosting theory: full literature review}
\label{sec:ghosting-lit-review}
Section~\ref{sec:ghosting-theory} gave the ghosting derivation in compressed form. This
section expands the literature review behind it in full, reproducing the
citation-fidelity discipline of the source document
(\texttt{Array\_Design/ghosting/THEORY.md}, verified read directly, mtime 2026-08-14
00:36) exactly: every claim below is labelled \textbf{[EXACT]} (a first-principles or
exact identity), \textbf{[PUBLISHED]} (a specific cited result applied as originally
derived), \textbf{[SYNTHESIS]} (a novel combination of citations not itself published as
a unified result), or \textbf{[UNTESTED]} (a claim requiring physical or psychophysical
confirmation this project has not obtained). This labelling convention is preserved
here deliberately, not smoothed into uniform prose, because it is itself part of this
project's own honesty discipline: a reader should be able to tell, claim by claim,
which statements in this section are load-bearing established physics and which are
this project's own proposed synthesis.

\subsection{Motivation: two real, sold products the original criterion could not explain}
An early criterion for this project treated adjacent holes at pitch $s$ as producing
retinal images separated by $\Delta r\approx s\,v_r/d_{\mathrm{gap}}$
\textbf{[EXACT, as originally labelled -- see below]}, and called a design ``ghost-free''
when $\Delta r$ was smaller than the single-hole blur radius $r_{EE}(P)$. Checked
against two real, commercially sold pinhole-eyewear patents -- KR~10-2081511 (pitch
3.2--4.8\,mm) and CN~201716487U (pitch 0.7--1.2\,mm), both with
$d_{\mathrm{gap}}\approx12$\,mm and $v_r\approx22$--$24$\,mm -- this criterion predicted
severe (30$\times$--400$\times$) violations for \emph{both} presumably-functional, sold
products. This direct contradiction between a theoretical prediction and two real
commercial devices is what motivated the full re-derivation this section and
Section~\ref{sec:ghosting-theory} document.

\begin{warnbox}{A label that did not survive re-derivation, kept visible rather than quietly corrected}
Equation~(1) above (the $\Delta r\approx s\,v_r/d_{\mathrm{gap}}$ form) was originally
marked \textbf{[EXACT]} in this project's own working notes. An independent
re-derivation subsequently found this label did not hold -- the formula is used nowhere
in this paper's reported results (Section~\ref{sec:ghosting-theory} uses the corrected
convolution-theorem derivation instead) -- and the source document's own convention is
to leave the original claim, and its original (incorrect) fidelity label, visible in
place with a note pointing to the correction, rather than silently edit history. This
paper follows that same convention here.
\end{warnbox}

\subsection{2.1 Lenslet-array merging theory (Gabor superlens)}
The closest existing precedent for a quantitative multi-channel merging condition in
the optics literature is Gabor-superlens theory (a microlens-array imaging
architecture, originating 1941, with modern quantitative treatments from 2014
\needsverify{exact citation for the specific 2014 treatment referenced in the source
document was not itself recorded there beyond the year -- not independently tracked
down for this paper}): mismatched, non-integer-ratio pitches in a two-array relay cause
imaging channels to, in the cited work's own language, ``unite to form superfocal
lines much smaller in number than the number of microlenses'' \textbf{[PUBLISHED]}.
This is cited here as evidence that the general class of problem -- multiple discrete
optical channels merging or failing to merge into a single perceived image -- is a
solved, quantitatively tractable one in an adjacent field, \emph{not} as a directly
transferable formula: a Gabor superlens is a two-array relay forming a real
intermediate image, physically a different system from a single perforated occluder
imaged directly by the eye's own optics, and no attempt is made in this project to
import the Gabor merging condition as-is.

\subsection{2.2 Direct evidence that multi-aperture ghosting is a real, reported clinical symptom}
The single most important literature finding underlying this section: Kim, Park, Park
\& Chun (2017, \emph{J Korean Med Sci} 32(5):850--857) \citep{kim2017} -- already
present in this project's own vetted citation library
(\texttt{stage4\_human\_vision/literature.py}, and independently used elsewhere in this
paper for a wholly separate purpose, Section~\ref{sec:external-validation}) -- directly
reports, in real human subjects wearing real multi-pinhole glasses, ``discomfort from
after-images such as honeycomb images, horizontal grating, and double or triple
vision'' \textbf{[PUBLISHED, direct quotation from the primary source]}, significantly
worse than for single-pinhole glasses tested in the same study. This finding overturns
the original framing of this project's own investigation: the question this project
initially set out to answer was ``why do real, sold multi-pinhole products show no
reported ghosting problem'' -- and the honest literature answer is that they
\emph{do} report exactly this problem, in a directly-quoted clinical source. The
question this section's derivation (Section~\ref{sec:ghosting-theory}) actually needed
to answer is therefore not ``why is ghosting absent,'' but ``why do the two specific
sold patents checked in Section~\ref{sec:ghosting-theory} not appear to address it
explicitly in their own documentation, and under exactly what geometric conditions does
the symptom become severe.''

\subsection{2.3 An independent real-world statement of the discrete-activation mechanism}
A separate, unrelated US patent (cited in the source document by its own title,
``Method for facilitating pinhole effect...''
\needsverify{full patent number not recorded in the source THEORY.md document beyond
this title fragment -- not independently tracked down for this paper}) states,
entirely independently of anything already in this project's own derivation: ``If two
pinholes are separated by less than the diameter of the pupil aperture, two pencils of
ray\ldots{} form two retinal images. Further, too large of a separation will result in
dead spots\ldots{} while too small\ldots{} will produce multiple images.''
\textbf{[PUBLISHED, direct quotation]}. This is significant specifically because it
states a \emph{discrete precondition} -- whether two or more holes are simultaneously
illuminating the pupil at all -- rather than a continuous margin comparison. This is
exactly the $N_{\mathrm{eff}}$ quantity this project's own active-pairs machinery
(\texttt{margins/active\_pairs.py}, Section~\ref{sec:three-regime}) already computes
via direct discrete-lattice enumeration, independently of this patent. The connection
between the patent's own stated mechanism and this project's own pre-existing
$N_{\mathrm{eff}}$ computation was not recognized until this literature review was
written -- i.e.\ this project's own Regime-A criterion ($N_{\mathrm{eff}}\le1$,
Equation~\ref{eq:three-regime}) turns out to be independently corroborated by real
patent-literature language describing the identical physical mechanism, discovered
after the fact rather than designed to match it.

\subsection{2.4 Periodic-artifact visibility: the screen-door-effect precedent}
Display engineering has a standard, mature treatment of an analogous problem: the
visibility of a display's own pixel/subpixel grid structure (the ``screen-door
effect''), which is assessed by expressing the grid's periodicity in cycles per
degree of visual angle, determining the resulting pattern's modulation depth, and
comparing that modulation depth against the human contrast sensitivity function's own
threshold at that spatial frequency \textbf{[PUBLISHED]}. This confirms that the
three-regime criterion's own frequency-domain framework
(Section~\ref{sec:ghosting-theory}, Equation~\ref{eq:three-regime}) is the same
established methodology display engineering already uses for a structurally analogous
periodic-artifact-visibility problem, not a bespoke construction invented for this
project alone. The same literature also confirms that a cruder, acuity-only rule (of
the same general type as this project's own original, later-superseded criterion,
Section~\ref{sec:ghosting-theory}) is a common, recognized first-pass approximation in
that adjacent field too -- i.e.\ this project's own original mistake was not a mistake
unique to this project, but a documented, recurring simplification in the broader
periodic-visual-artifact literature, corrected here by the same frequency-domain
methodology that field already uses to correct it.

\subsection{Synthesis: the three-regime criterion, and what remains untested}
Combining 2.2's discrete-activation precondition, 2.3's independent confirmation of the
same mechanism, and 2.4's frequency-domain visibility methodology, with this project's
own convolution-theorem derivation (Section~\ref{sec:ghosting-theory}), yields
Equation~\ref{eq:three-regime}'s three-regime criterion \textbf{[SYNTHESIS]} -- a novel
combination of these separately-published pieces, not itself a result published as a
unified criterion anywhere in the cited literature. The predicted danger-zone mechanism
(Regime C) is directionally consistent with Kim et al.'s directly-quoted clinical
symptom report (2.2 above), but this consistency is \emph{directional}, not a
quantitative validation: this project has not obtained psychophysical confirmation that
the specific Barten-CSF threshold crossing computed in
Section~\ref{sec:barten}--\ref{sec:harmonic} predicts \emph{which specific} real
designs a real wearer would report symptoms for, only that the general mechanism and
its general conclusion (dense real patent pitches are more dangerous than a
heuristic CSF model suggested) point the same direction as the one directly-quoted
clinical source available \textbf{[UNTESTED]} -- stated here explicitly as a limitation
of this specific sub-claim, in addition to the broader limitations already catalogued
in Chapter~\ref{ch:limitations}.
